\documentclass[12pt,a4paper]{article}
\usepackage[utf8]{inputenc}
\usepackage[turkish]{babel}
\usepackage{graphicx}
\usepackage{xcolor}
\usepackage{hyperref}
\usepackage{geometry}
\usepackage{titlesec}
\usepackage{fontawesome5} % İkonlar için
\usepackage{subcaption} % Yan yana görseller için ayar eklendi

% Sayfa kenar boşlukları
\geometry{a4paper, margin=1in}

% Marka Renkleri
\definecolor{primaryColor}{RGB}{59, 130, 246} % Mavi (Exploria Ana Renk)
\definecolor{secondaryColor}{RGB}{16, 185, 129} % Yeşil (Doğa/Keşif Rengi)
\definecolor{darkGray}{RGB}{55, 65, 81}

% Başlık Stilleri
\titleformat{\section}{\Large\bfseries\color{primaryColor}}{}{0em}{}[\titlerule]
\titleformat{\subsection}{\large\bfseries\color{secondaryColor}}{}{0em}{}

\title{\textbf{\Huge Exploria}\\[0.5em] \Large Dünyayı Oyun Alanınıza Çevirin}
\author{}
\date{}

\begin{document}

\maketitle

\begin{center}
    \textit{\textcolor{darkGray}{"Şehrin gizli kalmış hazinelerini keşfet, arkadaşlarınla takım ol ve kendi maceranı yaz!"}}
\end{center}

\vspace{0.5cm}

\section{\faMapMarkedAlt Yeni Nesil Keşif Deneyimi}
Exploria, geleneksel harita uygulamalarının çok ötesine geçerek şehir turlarını ve yerel maceraları \textbf{oyunlaştırılmış bir deneyime} dönüştürür. İster tarih kokan İstanbul/Fatih sokaklarında bir gezintiye çıkın, ister Ankara'nın kültürel mirasını harita üzerinden adımlayın; Exploria size "Gerçek dünyada bir oyun" hissi yaşatır.

\vspace{0.5cm}

\section{\faMagic Uygulama İçi Görünümler ve Temel Özellikler}

\subsection{1. Kişisel Profil ve Sosyal Ağ}
Kullanıcılar tamamen özelleştirilebilir bir profil ekranına sahiptir. Bu ekranda, bugüne kadar eklenen \textbf{arkadaş sayısı}, kazanılan \textbf{rozetler} ve kişisel bilgiler şık bir arayüzde sergilenir. Başarılarınızı anlık olarak bu ekrandan takip edebilir, profilinizi dilediğiniz gibi güncelleyebilirsiniz.

\begin{figure}[h!]
    \centering
    \includegraphics[width=0.45\textwidth]{profil.jpg}
    \caption{Modern ve dinamik profil yönetim ekranı.}
\end{figure}

\subsection{2. Renk Kodlu ve Kategorize Edilmiş Canlı Harita}
Exploria'nın kalbi olan harita ekranında, etrafınızdaki mekanlar (Cami, Tarihi Yapı, Çarşı \& Pazar vb.) renk kodlarıyla ayrılarak görsel bir şölen sunar. \textit{Sultanahmet Meydanı}, \textit{Yerebatan Sarnıcı} veya \textit{Topkapı Sarayı} gibi yüzlerce lokasyon özel ikonlarla haritada parlar.

Alt kısımda yer alan canlı sayaç \textit{(örn: Gezilen: 0 / 26 | 0/5925 hücre)} sayesinde, o bölgenin yüzde kaçını fethettiğinizi gerçek zamanlı olarak görebilir ve sisleri (Fog of War) adımladıkça yavaş yavaş dağıtabilirsiniz.

\begin{figure}[h!]
    \centering
    \includegraphics[width=0.45\textwidth]{harita.jpg}
    \caption{Kategorize edilmiş POI (Mekan) pinleri ve hücre keşif sistemi.}
\end{figure}

\newpage

\subsection{3. Çok Oyunculu (Co-op) Keşif Odaları}
Keşifleri tek başınıza yapmak zorunda değilsiniz! "İstanbul / Fatih Ekibi" gibi özel odalar kurarak arkadaşlarınızı davet edebilirsiniz. \textbf{Kurucu} yetkisiyle yönettiğiniz bu odalarda, takım arkadaşlarınızın katılımını bekleyebilir ve \textit{Keşfi Başlat} butonuyla haritadaki sisleri eş zamanlı, ortaklaşa dağıtarak mükemmel bir takım oyunu sergileyebilirsiniz.

\begin{figure}[h!]
    \centering
    \includegraphics[width=0.45\textwidth]{oda.jpg}
    \caption{Çok oyunculu (Multiplayer) takım odası bekleme salonu.}
\end{figure}

\subsection{4. Maceranız Kayıt Altında: Geçmiş Haritalar}
Gezdiğiniz, keşfettiğiniz ve fethettiğiniz haritalar asla kaybolmaz. \textbf{Geçmiş Haritalar} arşivi sayesinde, daha önce üzerinde yürüdüğünüz güzergahları, hangi tarihte o şehirde olduğunuzu kalıcı olarak profilinizde saklayabilir ve istediğiniz an anılarınızı tazeleyebilirsiniz.

\begin{figure}[h!]
    \centering
    \includegraphics[width=0.45\textwidth]{gecmis.jpg}
    \caption{Kullanıcının daha önce aktif ettiği haritaların arşiv görünümü.}
\end{figure}

\vspace{0.5cm}

\section{\faTrophy Neden Exploria?}
Sıradan bir yürüyüş aktivitesini, rekabet ve sosyalleşme duygusuyla harmanlıyoruz. Kullanıcı dostu şık arayüzü, gelişmiş harita renklendirmesi ve güçlü "Çok Oyunculu Ekip" altyapısıyla Exploria; sokağa çıkmak için ihtiyaç duyduğunuz o ekstra motivasyonun ta kendisidir!

\vfill
\begin{center}
    \textbf{Exploria} - \textit{Dışarı Çık ve Keşfet!}
\end{center}

\end{document}
